%% ch03.tex — Chapter 3: Distributed File Systems: GFS and HDFS
%% Part of "Data Engineering" textbook

\chapter{Distributed File Systems: GFS and HDFS}
\label{ch:hdfs}

\epigraph{\textit{``We have designed the system for a usage pattern in which
most files are mutated by appending new data rather than overwriting
existing data. Random writes within a file are practically non-existent.''}}{Ghemawat, Gobioff \& Leung, SOSP 2003}

\vspace{4pt}

\begin{chapterlearning}
After completing this chapter, you will be able to:
\begin{enumerate}
  \item Explain why conventional file systems (NFS, POSIX, SAN) cannot
        meet the storage requirements of petabyte-scale data engineering,
        and identify the specific assumptions they violate.
  \item Describe the GFS architecture---master, chunkservers, clients,
        chunk size, and metadata design---and explain the key design
        decisions that made it suitable for Google's workload.
  \item Draw and explain the HDFS architecture: NameNode, DataNodes,
        Secondary NameNode, and their interactions.
  \item Trace the complete HDFS read and write paths step by step,
        including pipeline replication, checksum verification, and
        lease management.
  \item Explain rack-aware replication: the default three-replica
        placement policy and why it balances fault tolerance against
        write overhead.
  \item Describe the NameNode high-availability (HA) architecture with
        Active/Standby nodes, ZooKeeper leader election, and shared
        Journal Nodes.
  \item Identify HDFS's design limitations---the small file problem,
        no random writes, low-latency access---and contrast HDFS
        with cloud object stores (Amazon S3, Google Cloud Storage).
\end{enumerate}
\end{chapterlearning}

\bigskip

Chapter~\ref{ch:integration} established how data moves between systems
and how heterogeneous sources are unified. Before any of that
integration logic can execute, however, the raw data must live
somewhere. In 2003, Google faced a storage problem that no existing
file system could solve: billions of web pages, totalling petabytes of
data, needed to be stored reliably across thousands of commodity
servers, indexed sequentially at enormous throughput, and kept
available even as individual servers failed---which, at that scale,
happened multiple times per day.

The solution Google published was the \term{Google File System} (GFS)
\parencite{GFS2003}. Three years later, the Apache Hadoop project
reimplemented and open-sourced the same design as the \term{Hadoop
Distributed File System} (HDFS). HDFS became the foundational storage
layer for an entire ecosystem of distributed data processing tools---
MapReduce, Hive, Pig, HBase, and eventually Spark---and it remains in
active production use at organisations ranging from Yahoo! and Facebook
to LinkedIn and the New York Times. Understanding HDFS is understanding
the bedrock on which the rest of this book's processing stack rests.

%%====================================================================
\section{The Storage Problem at Petabyte Scale}
\label{sec:storage-problem}
%%====================================================================

A \term{file system} is the software layer responsible for organising,
naming, accessing, and protecting data stored on physical storage
devices. The file system abstracts away the physical device---magnetic
platters, solid-state cells, or optical discs---and presents the
programmer with a uniform interface of named files in a hierarchical
directory structure. For four decades, this model served the computing
industry well. It continues to serve it well for most workloads today.

The model breaks at petabyte scale for three fundamental reasons.

\textbf{Single-machine capacity.} A conventional file system is
managed by a single host. Even the largest shared file systems---
Network File System (NFS) and SAN-attached volumes---ultimately
depend on a finite number of storage controllers. By 2003, Google's
web crawl was producing hundreds of terabytes per month, a rate that
had already surpassed what any single storage system could absorb.
The only solution was to distribute data across a cluster of
independent machines.

\textbf{Throughput requirements.} Web indexing requires reading an
entire dataset sequentially to reprocess it---there is no practical
way to index a subset of the web. A single high-performance SCSI disk
in 2003 delivered approximately 50 MB/s of sequential read throughput.
Reading one petabyte from a single disk would take over 230 days.
Distributing the data across 1,000 disks and reading them in parallel
reduces this to approximately five hours---a feasible engineering
timeline.

\textbf{Hardware failure as normal operation.} A conventional file
system---and the RDBMS systems that sit on top of it---treats hardware
failure as an exceptional condition to be recovered from. At Google's
scale in 2003, a cluster of a few thousand commodity PCs experienced
approximately three node failures per day as normal operation: disk
bearing failures, memory errors, motherboard failures, and network card
failures, all occurring on a predictable statistical schedule
\parencite{GFS2003}. Any storage system that treated such failures as
exceptional events would be unavailable for recovery most of the time.
The failure model must instead treat hardware failure as \emph{the
expected case} and design accordingly.

\begin{table}[ht]
\centering
\caption{Comparison of storage architectures for big data workloads}
\label{tab:storage-comparison}
\renewcommand{\arraystretch}{1.4}
\begin{tabular}{L{2.0cm} L{2.4cm} L{2.4cm} L{4.0cm}}
\toprule
\textbf{System} & \textbf{Max capacity} & \textbf{Fault model} & \textbf{Reason it fails at big data scale} \\
\midrule
POSIX local FS  & Single disk / array & Single point of failure & No horizontal scaling; one server limit \\
NFS             & NAS head + array    & NAS head is SPOF       & Centralized metadata and I/O bottleneck \\
SAN             & Storage array       & Expensive RAID         & Proprietary hardware; prohibitive cost/PB \\
RDBMS blob      & Server RAM + disk   & WAL recovery           & Tuned for random access; not sequential scan \\
\textbf{GFS / HDFS} & \textbf{Cluster-wide} & \textbf{Expects failure} & \textbf{Designed for this workload} \\
\bottomrule
\end{tabular}
\end{table}

%%====================================================================
\section{The Google File System}
\label{sec:gfs}
%%====================================================================

The Google File System, published by Ghemawat, Gobioff, and Leung at
the ACM Symposium on Operating Systems Principles (SOSP) in October
2003, is one of the most consequential systems papers of the past two
decades \parencite{GFS2003}. It describes a purpose-built distributed
file system designed not to satisfy general-purpose file system
requirements---POSIX compliance, low-latency random access, small
file support---but to satisfy the specific requirements of Google's
production workloads.

%--------------------------------------------------------------------
\subsection{Workload Assumptions and Design Decisions}
\label{subsec:gfs-workload}
%--------------------------------------------------------------------

GFS was designed around a careful empirical study of Google's actual
workloads. The authors observed that virtually all files in Google's
production environment were large---multiple gigabytes to hundreds of
gigabytes---and that data access patterns were dominated by sequential
reads and sequential appends, with random writes essentially absent.
Files were written once, read many times, and never modified in place
after initial creation. Concurrent reads from many clients were the
norm; concurrent random writes to the same file were not.

These observations drove six key design decisions that distinguish GFS
from conventional file systems:

\begin{enumerate}
  \item \textbf{Large chunk size (64 MB).} Files are divided into
        fixed-size \term{chunks} of 64\MB. This is orders of magnitude
        larger than the 4-\KB\ blocks used by conventional file systems.
        Large chunks reduce the amount of metadata the master must
        maintain (one entry per chunk rather than one per block), reduce
        client-to-master communication (a client needs the location of
        fewer chunks to read a large file), and encourage direct TCP
        connections between clients and chunkservers for sustained
        sequential I/O.

  \item \textbf{Centralised, in-memory metadata.} The GFS master holds
        all file system metadata---the file namespace, chunk-to-file
        mappings, chunk locations, access control lists---entirely in
        RAM. This enables sub-millisecond metadata lookups. The master
        has a single-machine capacity for metadata because large chunk
        sizes keep the number of chunks manageable: at 64\MB\ per chunk,
        one petabyte of data requires only $\sim$16 million chunk
        entries, which fits comfortably in modern server RAM.

  \item \textbf{No data caching.} GFS explicitly does not cache file
        data at either clients or chunkservers. The workload is
        dominated by streaming reads that iterate through files once;
        cached data would be evicted before it could be reused. The
        kernel's buffer cache handles any incidental temporal locality.

  \item \textbf{Write-by-append semantics.} GFS optimises for atomic
        append operations rather than random writes. A client can append
        data to any file at the current end-of-file position; this
        corresponds to how log files, web crawl outputs, and search
        index segments are naturally produced.

  \item \textbf{Fault tolerance through replication, not redundant
        hardware.} Each 64\MB\ chunk is stored on three chunkservers
        by default. If one fails, the master re-replicates the
        affected chunks on a surviving server. No RAID arrays or
        specialised storage hardware are required.

  \item \textbf{Relaxed consistency.} GFS provides a consistency model
        weaker than POSIX. Concurrent writes to the same region from
        different clients may produce an interleaved, ``defined but
        inconsistent'' result. Google's applications were designed to
        tolerate this---for example, by writing unique chunk identifiers
        with each record and using deduplication at read time.
\end{enumerate}

%--------------------------------------------------------------------
\subsection{GFS Architecture}
\label{subsec:gfs-arch}
%--------------------------------------------------------------------

A GFS cluster consists of three types of components: one \term{GFS
master}, multiple \term{chunkservers}, and multiple \term{clients}.

The \textbf{GFS master} maintains all file system metadata in RAM. It
does not participate in data I/O between clients and chunkservers; once
a client has obtained the chunk location from the master, all data
transfer occurs directly between the client and the chunkserver. This
separation of the control plane (master) from the data plane
(chunkservers and clients) is the central architectural decision that
enables GFS to scale: the master handles only lightweight metadata
operations, while the bulk data flows on separate direct connections.

The master communicates with chunkservers through two periodic
mechanisms. Each chunkserver sends a \term{heartbeat} to the master
every few seconds to confirm it is operational; the absence of
heartbeats for a configurable period causes the master to mark the
chunkserver as failed. Each chunkserver also sends a periodic
\term{chunk report} listing all the chunk replicas it currently holds,
which allows the master to reconstruct its view of chunk locations
after a restart.

\textbf{Chunkservers} store chunks as plain files on their local Linux
file systems. They do not cache chunks in memory; the kernel's buffer
cache provides whatever temporal locality exists. Each chunkserver is
responsible for maintaining checksum metadata for every chunk it stores,
and it verifies checksums on every read to detect data corruption.

\textbf{Clients} implement the GFS API (which is not POSIX-compatible).
A client contacts the master to resolve a file name to a list of chunk
handles and their chunkserver locations, then communicates directly
with the appropriate chunkservers for data I/O. Clients cache the
chunk location information locally for a limited time, reducing
master load.

\begin{researchspotlight}{Ghemawat, Gobioff \& Leung (2003) --- The Google File System}
\textbf{Citation:} Ghemawat, S., Gobioff, H., \& Leung, S.-T.\ (2003).
The Google File System. \textit{Proceedings of the 19th ACM Symposium
on Operating Systems Principles (SOSP)}, pp.\ 29--43.
\parencite{GFS2003}

\smallskip\noindent
\textbf{Key insight:} The authors' central argument is that
\emph{workload-specific design} outperforms general-purpose design at
extreme scale. By deliberately abandoning POSIX compliance, random
write support, and fine-grained access control, GFS achieved
throughput and fault tolerance impossible with any off-the-shelf
file system.

\smallskip\noindent
\textbf{Technical contributions:}
\begin{itemize}
  \item Demonstrated that a centralised metadata server (the master)
        is not a scalability bottleneck when (a) metadata is kept
        entirely in RAM, (b) chunk sizes are large enough that the
        total metadata volume is manageable, and (c) data I/O bypasses
        the master entirely.
  \item Introduced \emph{atomic record append}---an operation that
        allows multiple clients to append to the same file
        concurrently without prior coordination, enabling
        producer-consumer pipelines over shared log files.
  \item Quantified failure rates in commodity clusters and provided
        the engineering argument that replication (software-based
        redundancy on commodity hardware) is more cost-effective than
        hardware RAID at petabyte scale.
\end{itemize}

\smallskip\noindent
\textbf{Impact today:} GFS is cited in over 14,000 academic papers---
one of the most-cited systems papers ever published. It directly
inspired HDFS (the open-source reimplementation), Amazon S3, Microsoft
Azure Blob Storage, and Google's own successor system, Colossus.
The paper is assigned reading at Stanford CS246, MIT 6.830, and CMU
15-721.
\end{researchspotlight}

%%====================================================================
\section{HDFS Architecture}
\label{sec:hdfs-arch}
%%====================================================================

The Hadoop Distributed File System is the open-source Apache
implementation of the GFS design. It was initially developed at Yahoo!
by Doug Cutting and Mike Cafarella---who had been building a web
crawler called Nutch and needed a production-ready open-source
replacement for GFS---and contributed to the Apache Software Foundation
in 2006. HDFS translates the GFS architecture into Java, with some
naming changes (GFS master $\to$ NameNode, chunkserver $\to$ DataNode,
chunk $\to$ block) and some incremental design improvements.

%--------------------------------------------------------------------
\subsection{The NameNode}
\label{subsec:namenode}
%--------------------------------------------------------------------

The \term{NameNode} is HDFS's master process. It is the centralised
metadata server responsible for maintaining the file system namespace
and the block-to-DataNode mapping. Like the GFS master, it holds all
metadata entirely in RAM; the typical NameNode memory footprint is
approximately 150 bytes per file and 150 bytes per block, meaning that
one million files average $\sim$300\MB\ of heap---manageable on a
modern server.

The NameNode's metadata is persisted to disk in two complementary
structures. The \term{FsImage} is a complete, point-in-time snapshot
of the file system namespace---a serialised checkpoint of all metadata.
The \term{EditLog} is a write-ahead log of all namespace mutations
since the last checkpoint: file creations, deletions, permission
changes, and block allocations. On startup, the NameNode replays the
EditLog against the last FsImage to reconstruct the current namespace
state. Under normal operation, the EditLog grows continuously; it is
periodically merged into the FsImage in a process called
\term{checkpointing}.

The NameNode does \emph{not} store the physical locations of blocks in
persistent metadata. Instead, block locations are rebuilt from scratch
at startup through the DataNode \term{block report}---each DataNode
reports the full list of blocks it currently holds when it connects to
the NameNode. This design avoids the risk of stale location data in
the FsImage becoming inconsistent with the actual DataNode state after
failures or maintenance.

\begin{defbox}{NameNode Responsibilities}
\begin{enumerate}
  \item \textbf{Namespace management:} maintain the directory tree and
        file-to-block mapping.
  \item \textbf{Block management:} track which DataNodes hold each
        block replica; initiate re-replication when a replica is lost.
  \item \textbf{Access control:} enforce file permissions on all client
        requests.
  \item \textbf{Lease management:} grant and revoke write leases that
        determine which client may currently write to a file.
  \item \textbf{Load balancing:} coordinate the HDFS Balancer, which
        moves blocks between DataNodes to equalise disk utilisation.
\end{enumerate}
\end{defbox}

%--------------------------------------------------------------------
\subsection{DataNodes}
\label{subsec:datanodes}
%--------------------------------------------------------------------

\term{DataNodes} are the worker processes that store actual block data.
Each DataNode manages its local storage---one or more hard disks or
SSDs---and serves block read and write requests directly from clients.
DataNodes are designed to run on commodity hardware; a typical DataNode
configuration in 2024 uses 12--24 spinning hard disks of 8--18\TB\
each, providing 96--432\TB\ of raw capacity per node.

DataNodes communicate with the NameNode through two periodic messages.
A \term{heartbeat} (sent every 3 seconds by default) confirms the
DataNode is alive and operational. A \term{block report} (sent every 6
hours by default, or on demand at startup) lists all block replicas the
DataNode currently holds. The NameNode uses both to maintain its live
view of the cluster's block map. If a heartbeat is absent for more than
10.5 minutes (the default \texttt{dfs.heartbeat.interval} $\times$
\texttt{dfs.namenode.heartbeat.recheck-interval}), the NameNode marks
the DataNode as dead and schedules re-replication of all blocks it held.

DataNodes also perform \term{block scanning}: each DataNode
independently verifies the checksums of all its stored blocks on a
configurable schedule (default: every three weeks). A block that fails
checksum verification is reported to the NameNode as corrupted, and
the NameNode responds by scheduling re-replication from a healthy
replica.

%--------------------------------------------------------------------
\subsection{The Secondary NameNode}
\label{subsec:secondary-nn}
%--------------------------------------------------------------------

\begin{caution}{The Secondary NameNode is Not a Hot Standby}
Despite its misleading name, the \textbf{Secondary NameNode} is not a
standby that takes over if the primary NameNode fails. It performs a
single function: it periodically merges the EditLog into the FsImage,
producing a new, updated FsImage that it sends back to the primary
NameNode. This \emph{checkpointing} service reduces the EditLog size
and therefore the NameNode startup time after a crash. If the primary
NameNode fails, the Secondary NameNode's checkpoint can be used to
recover the namespace, but this requires a manual failover process
and results in some minutes of downtime. True hot standby failover
requires the High Availability architecture described in
Section~\ref{sec:hdfs-ha}.
\end{caution}

The checkpointing cycle works as follows. The Secondary NameNode
fetches the current FsImage and EditLog from the primary, applies all
EditLog transactions to the FsImage in memory, writes the resulting
merged FsImage to disk, and sends it back to the primary. The primary
then begins writing subsequent namespace mutations to a fresh, empty
EditLog. The net result is that the EditLog never grows indefinitely;
checkpointing occurs every hour by default, or when the EditLog exceeds
one million transactions.

%--------------------------------------------------------------------
\subsection{Block Storage}
\label{subsec:blocks}
%--------------------------------------------------------------------

\term{Blocks} are the fundamental unit of storage in HDFS. The default
block size is 128\MB\ (increased from the original 64\MB\ in Hadoop
2.x). When a client writes a file, HDFS divides it into 128\MB\ blocks
and distributes those blocks across DataNodes; each block is stored
three times by default (the \term{replication factor}).

The large block size---32,768 times larger than a typical ext4 file
system block---serves several purposes. It reduces the number of
metadata entries the NameNode must maintain. It enables sustained
sequential I/O at full disk throughput without per-block overhead. And
it reduces the number of round trips between client and NameNode
required to locate all blocks of a large file.

A notable edge case: the last block of a file may be smaller than
128\MB. HDFS does not pad it to 128\MB; the block occupies only as
much disk space as the actual data. A 1-\GB\ file consists of eight
full 128\MB\ blocks, not nine.

%--------------------------------------------------------------------
\subsection*{The Complete HDFS Cluster Architecture}
\label{subsec:hdfs-arch-overview}
%--------------------------------------------------------------------

Figure~\ref{fig:hdfs-arch} assembles all three components into a single
view of the HDFS cluster architecture. The key point to internalise is
the separation of the \emph{control plane} (all metadata operations that
involve the NameNode) from the \emph{data plane} (all block transfers
that flow directly between clients and DataNodes). The NameNode processes
millions of metadata RPC calls per day but never touches a byte of block
data; every byte of data is transferred directly between clients and
DataNodes on independent TCP connections.

\begin{figure}[ht]
\centering
\begin{tikzpicture}[
    nn/.style={rectangle, draw=DEBlue, fill=DELightBlue!50, rounded corners=5pt,
               minimum width=8.2cm, minimum height=2.1cm},
    dn/.style={rectangle, draw=DEGreen!70!black, fill=DELightGreen, rounded corners=4pt,
               minimum width=2.2cm, minimum height=1.8cm, align=center,
               font=\small\sffamily},
    snn/.style={rectangle, draw=DEGray!65, fill=DEGray!10, rounded corners=4pt,
                minimum width=2.4cm, minimum height=1.5cm, align=center,
                font=\small\sffamily},
    cl/.style={rectangle, draw=DEAccent, fill=DELightAccent, rounded corners=4pt,
               minimum width=2.0cm, minimum height=0.85cm, align=center,
               font=\small\bfseries\sffamily},
    ibox/.style={rectangle, draw=DEBlue!40, fill=white, rounded corners=2pt,
                 minimum width=1.7cm, minimum height=0.6cm,
                 font=\tiny\sffamily, align=center},
    arr/.style={-{Stealth[length=4pt,width=4pt]}, thick},
    darr/.style={{Stealth[length=4pt,width=4pt]}-{Stealth[length=4pt,width=4pt]}, thick},
    lbl/.style={font=\scriptsize\sffamily, align=center}
  ]

  %--- NameNode ---
  \node[nn] (nn) at (4.8, 6.0) {};
  \node[font=\small\bfseries\sffamily\color{DEBlue}] at (4.8, 7.0)
       {NameNode \textnormal{(Master)}};
  \node[ibox] at (2.7, 6.05) {File Namespace\\(in RAM)};
  \node[ibox] at (4.8, 6.05) {FsImage\\(disk)};
  \node[ibox] at (6.9, 6.05) {EditLog\\(disk)};
  \node[font=\tiny\sffamily, text=DEGray] at (4.8, 5.28)
       {Metadata plane only --- no block data passes through};

  %--- Secondary NameNode ---
  \node[snn] (snn) at (11.5, 6.2) {Secondary\\NameNode\\[2pt]%
       {\tiny\textit{(checkpointing)}}};

  %--- DataNodes ---
  \node[dn] (dn1) at (1.6, 2.0)
       {DataNode 1\\[3pt]{\scriptsize Rack A}\\[4pt]{\tiny blocks on disk}};
  \node[dn] (dn2) at (4.8, 2.0)
       {DataNode 2\\[3pt]{\scriptsize Rack B}\\[4pt]{\tiny blocks on disk}};
  \node[dn] (dn3) at (8.0, 2.0)
       {DataNode 3\\[3pt]{\scriptsize Rack B}\\[4pt]{\tiny blocks on disk}};

  %--- Client ---
  \node[cl] (cl) at (11.5, 2.5) {Client};

  %--- Heartbeat + Block Report: DataNode <-> NameNode ---
  % NN bottom is at y = 6.0 - 1.05 = 4.95
  % DN top  is at y = 2.0 + 0.90 = 2.90
  \draw[darr, DEGreen!65!black]
    (dn1.north) -- (1.6, 4.95)
    node[midway, left, lbl, text=DEGreen!65!black]
         {\tiny heartbeat\\[-2pt]\tiny block report};
  \draw[darr, DEGreen!65!black]
    (dn2.north) -- (4.8, 4.95);
  \draw[darr, DEGreen!65!black]
    (dn3.north) -- (8.0, 4.95)
    node[midway, right, lbl, text=DEGreen!65!black]
         {\tiny heartbeat\\[-2pt]\tiny block report};

  %--- Client <-> NameNode: metadata RPC (control plane, step 1) ---
  % nn.east is at x = 4.8 + 4.1 = 8.9, y = 6.0
  \draw[darr, DEAccent]
    (cl.west) to[out=170, in=0]
    node[midway, above, lbl, text=DEAccent]
         {\tiny\textbf{(1)} metadata RPC / block locations}
    (nn.east);

  %--- Client -> DataNode: direct data I/O (data plane, step 2) ---
  \draw[arr, DEBlue, dashed, thick]
    (cl.south) to[out=220, in=-10]
    node[midway, below, lbl, text=DEBlue]
         {\tiny\textbf{(2)} data I/O (TCP direct)}
    (dn3.east);

  %--- NameNode -> Secondary NameNode: send FsImage + EditLog ---
  \draw[arr, DEGray!60]
    (nn.north east) to[out=25, in=155]
    node[midway, above, lbl, text=DEGray!60]
         {\tiny FsImage + EditLog}
    (snn.west);

  %--- Secondary NameNode -> NameNode: return merged FsImage ---
  \draw[arr, DEGray!60]
    (snn.south) to[out=255, in=5]
    node[pos=0.45, right, lbl, text=DEGray!60]
         {\tiny merged FsImage\\[-2pt]\tiny (checkpoint)}
    (nn.east);

  %--- Rack dashed boundaries ---
  \draw[dashed, DEGray!40, rounded corners=5pt]
       (0.4, 0.8) rectangle (2.8, 3.2);
  \node[font=\tiny\sffamily, text=DEGray!55] at (1.6, 0.5) {Rack A};

  \draw[dashed, DEGray!40, rounded corners=5pt]
       (3.6, 0.8) rectangle (9.2, 3.2);
  \node[font=\tiny\sffamily, text=DEGray!55] at (6.4, 0.5) {Rack B};

\end{tikzpicture}
\caption{Complete HDFS cluster architecture. The NameNode (master) stores
         all file system metadata in RAM (file namespace, block-to-DataNode
         map) and persists mutations to the EditLog and FsImage. DataNodes
         (workers) store blocks on local disks and communicate with the
         NameNode via periodic heartbeats and block reports. Clients obtain
         block locations from the NameNode \emph{(step~1 --- control plane)},
         then read/write block data directly to DataNodes \emph{(step~2 ---
         data plane)}, bypassing the NameNode entirely. The Secondary
         NameNode checkpoints the EditLog into the FsImage to bound restart
         time after a NameNode crash.}
\label{fig:hdfs-arch}
\end{figure}

%%====================================================================
\section{Replication and Fault Tolerance}
\label{sec:replication}
%%====================================================================

HDFS achieves fault tolerance through \term{block replication}: storing
each block on multiple DataNodes such that if any individual DataNode
fails, the remaining replicas ensure data availability. The default
replication factor of 3 means that the failure of any single DataNode
does not result in any data loss or unavailability.

%--------------------------------------------------------------------
\subsection{Rack-Aware Replication}
\label{subsec:rack-aware}
%--------------------------------------------------------------------

Naive replication---placing all three replicas on randomly chosen
DataNodes---would provide poor fault tolerance in a real data centre.
Physical servers in a data centre are organised in \term{racks}: sets
of 20--40 servers sharing a top-of-rack switch. A rack-level failure
(a switch failure or a power distribution unit trip) takes down every
server in the rack simultaneously. If all three replicas of a block
happened to reside on nodes within the same rack, a rack failure would
destroy all three replicas.

HDFS addresses this with \term{rack-aware replication}. The default
policy for replication factor 3 places:
\begin{itemize}
  \item \textbf{Replica 1}: on the same node as the writing client (or
        a random node if the client is not running on a DataNode).
  \item \textbf{Replica 2}: on a DataNode in a \emph{different} rack
        from Replica 1.
  \item \textbf{Replica 3}: on a \emph{different} DataNode in the same
        rack as Replica 2.
\end{itemize}

This placement balances three competing objectives. Placing Replica 1
locally (or on a nearby node) minimises write latency for the initial
block placement. Placing Replica 2 on a different rack ensures that a
single rack failure cannot destroy all replicas. Placing Replica 3 in
the same rack as Replica 2 (rather than a third rack) reduces cross-rack
network traffic: reads can often be satisfied from within the same rack
as the client, reducing the load on the inter-rack network bandwidth.

\begin{figure}[ht]
\centering
\begin{tikzpicture}[
    server/.style={rectangle, draw=DEBlue!50, fill=DELightBlue,
                   minimum width=1.4cm, minimum height=0.6cm,
                   font=\scriptsize\sffamily, align=center},
    rack/.style={rectangle, draw=DEGray!60, fill=DEGray!8,
                 rounded corners=3pt, inner sep=8pt},
    sw/.style={rectangle, draw=DEGray!60, fill=DEGray!25,
               minimum width=3.2cm, minimum height=0.4cm,
               font=\scriptsize\sffamily},
    wire/.style={draw=DEGray!50, thick},
    replica/.style={circle, minimum size=0.28cm, inner sep=0pt,
                    font=\scriptsize\bfseries, text=white}
  ]
  % Rack A
  \begin{scope}[local bounding box=rackA]
    \node[server] (n1) at (0, 3.6) {DN-1\\{\tiny\itshape Replica 1}};
    \node[server] (n2) at (0, 2.7) {DN-2};
    \node[server] (n3) at (0, 1.8) {DN-3};
    \node[sw]     (swA) at (0, 1.2) {Switch A};
    \node[rack, fit=(n1)(n2)(n3)(swA),
          label={[font=\small\bfseries\sffamily\color{DEBlue}]above:Rack A}] {};
  \end{scope}

  % Rack B
  \begin{scope}[xshift=4.5cm, local bounding box=rackB]
    \node[server] (n4) at (0, 3.6) {DN-4\\{\tiny\itshape Replica 2}};
    \node[server] (n5) at (0, 2.7) {DN-5\\{\tiny\itshape Replica 3}};
    \node[server] (n6) at (0, 1.8) {DN-6};
    \node[sw]     (swB) at (0, 1.2) {Switch B};
    \node[rack, fit=(n4)(n5)(n6)(swB),
          label={[font=\small\bfseries\sffamily\color{DEBlue}]above:Rack B}] {};
  \end{scope}

  % Inter-rack switch
  \node[sw, minimum width=4cm] (core) at (2.25, 0.3)
        {Core / Aggregation Switch};

  \draw[wire] (swA) -- (core);
  \draw[wire] (swB) -- (core);

  % Replica indicators
  \node[replica, fill=DERed]    at (0.9, 3.6) {1};
  \node[replica, fill=DEAccent] at (0.9, 3.6+0) {};  % already labelled
  \node[replica, fill=DEGreen]  at (4.5+0.9, 3.6) {2};
  \node[replica, fill=DEBlue]   at (4.5+0.9, 2.7) {3};

  % Legend
  \node[font=\scriptsize\sffamily, align=left] at (9.5, 2.8) {%
    \textbf{Replica placement}\\[3pt]
    R1: local write node (DN-1)\\
    R2: different rack (DN-4)\\
    R3: same rack as R2 (DN-5)\\[5pt]
    \textit{A single rack failure}\\
    \textit{cannot lose all replicas.}};
\end{tikzpicture}
\caption{HDFS rack-aware replication with replication factor 3.
         Two replicas in Rack~B; one in Rack~A. A Rack~A failure
         leaves Replicas~2 and~3 available.}
\label{fig:rack-aware}
\end{figure}

%--------------------------------------------------------------------
\subsection{Failure Detection and Re-replication}
\label{subsec:re-replication}
%--------------------------------------------------------------------

When the NameNode detects that a DataNode is dead (missed heartbeats
for more than the configured threshold), it immediately identifies
every block for which that DataNode held a replica. For each such
block, the NameNode checks whether the remaining replicas satisfy the
replication factor. If a block now has fewer than the required number
of replicas, the NameNode schedules \term{re-replication}: it instructs
one of the surviving DataNodes that holds a copy to transfer the block
to a new DataNode, restoring the full replication factor.

Re-replication is prioritised by urgency. Blocks that have only one
surviving replica (single-replica blocks) are re-replicated first,
since they are at immediate risk of data loss if the last surviving
DataNode also fails. Blocks with two surviving replicas are re-replicated
next. The re-replication process consumes cluster network bandwidth; in
a large cluster where many DataNodes are replaced simultaneously (e.g.,
during a scheduled maintenance window), the NameNode throttles
re-replication to avoid saturating the network.

%%====================================================================
\section{HDFS Read and Write Paths}
\label{sec:hdfs-io}
%%====================================================================

Understanding how data actually flows in and out of HDFS---not merely
the architecture of the components but the sequence of operations---is
essential for diagnosing performance problems and for understanding why
HDFS makes the consistency trade-offs it does.

%--------------------------------------------------------------------
\subsection{The Read Path}
\label{subsec:read-path}
%--------------------------------------------------------------------

Reading a file from HDFS involves the following sequence:

\begin{enumerate}
  \item \textbf{Open.} The client calls \texttt{FileSystem.open()}.
        The HDFS client library contacts the NameNode via RPC and
        requests the block locations for the first few blocks of the
        target file.

  \item \textbf{Block location response.} The NameNode returns an
        ordered list of DataNode addresses for each requested block,
        sorted by their network proximity to the client. The list
        prioritises DataNodes on the same node as the client, then
        the same rack, then other racks---the data locality principle
        from Chapter~\ref{ch:foundations}.

  \item \textbf{Direct DataNode connection.} The client connects
        directly to the closest DataNode for the first block and
        streams the data. The client reads data from the DataNode
        through a TCP socket; no data passes through the NameNode.

  \item \textbf{Checksum verification.} As each 512-byte
        \term{checksum chunk} is read, the client verifies its CRC-32C
        checksum against the stored checksum. A mismatch indicates
        data corruption; the client reports the bad block to the
        NameNode and retries on an alternative replica.

  \item \textbf{Block boundary.} When one block is fully read, the
        client contacts the NameNode for the locations of subsequent
        blocks (if not already cached) and opens a new connection to
        the appropriate DataNode.

  \item \textbf{Close.} When all blocks have been read, the input
        stream is closed.
\end{enumerate}

\begin{figure}[ht]
\centering
\begin{tikzpicture}[
    actor/.style={rectangle, draw=DEBlue!60, fill=DELightBlue!40, rounded corners=4pt,
                  minimum width=2.8cm, minimum height=0.75cm, align=center,
                  font=\small\bfseries\sffamily},
    step/.style={rectangle, draw=DEGreen!60, fill=DELightGreen!60, rounded corners=3pt,
                 minimum width=3.4cm, minimum height=0.6cm, align=center,
                 font=\scriptsize\sffamily},
    arr/.style={-{Stealth[length=4pt,width=4pt]}, thick},
    lbl/.style={font=\scriptsize\sffamily, align=center}
  ]

  % Life-line headers
  \node[actor] (c) at (0,   8.5) {Client};
  \node[actor] (n) at (5.5, 8.5) {NameNode};
  \node[actor] (d) at (11,  8.5) {DataNode};

  % Life-lines (dashed vertical)
  \draw[dashed, DEGray!40] (0,   8.1) -- (0,   0.0);
  \draw[dashed, DEGray!40] (5.5, 8.1) -- (5.5, 0.0);
  \draw[dashed, DEGray!40] (11,  8.1) -- (11,  0.0);

  % Step labels (left margin)
  \foreach \y/\txt in {
    7.7/{\textbf{1.}\ open()},
    6.7/{\textbf{2.}\ block locations (sorted by proximity)},
    5.5/{\textbf{3.}\ TCP connect to closest DataNode},
    4.3/{\textbf{4.}\ stream block data (512-byte chunks)},
    3.2/{\textbf{5.}\ CRC checksum verification},
    2.1/{\textbf{6.}\ next block (repeat 3--5)},
    0.9/{\textbf{7.}\ close()}}
  {
    \node[lbl, text=DEGray, anchor=east] at (-0.3, \y) {};
  }

  % 1. Client -> NameNode: open()
  \draw[arr, DEAccent]
    (0, 7.7) -- (5.3, 7.7)
    node[midway, above, lbl, text=DEAccent] {\tiny open()};

  % 2. NameNode -> Client: block locations
  \draw[arr, DEAccent]
    (5.3, 6.7) -- (0.2, 6.7)
    node[midway, above, lbl, text=DEAccent]
         {\tiny block locations (ordered by network proximity)};

  % 3. Client -> DataNode: TCP connect
  \draw[arr, DEBlue]
    (0.2, 5.5) -- (10.7, 5.5)
    node[midway, above, lbl, text=DEBlue]
         {\tiny TCP connect to closest DN};

  % 4. DataNode -> Client: stream data
  \draw[arr, DEGreen!70!black, very thick]
    (10.7, 4.3) -- (0.2, 4.3)
    node[midway, above, lbl, text=DEGreen!70!black]
         {\tiny stream block data at full disk bandwidth};

  % 5. Client: checksum verify (self-loop)
  \draw[arr, DEBlue!60]
    (0.2, 3.5) to[out=30, in=-30] (0.2, 3.0)
    node[pos=0.5, right, lbl, text=DEBlue!60]
         {\tiny CRC-32C verify};

  % 6. Client -> DataNode: next block
  \draw[arr, DEBlue]
    (0.2, 2.1) -- (10.7, 2.1)
    node[midway, above, lbl, text=DEBlue]
         {\tiny open next block (NN re-consulted if needed)};

  \draw[arr, DEGreen!70!black, very thick]
    (10.7, 1.6) -- (0.2, 1.6)
    node[midway, above, lbl, text=DEGreen!70!black]
         {\tiny stream next block};

  % 7. Client -> NameNode: close
  \draw[arr, DEAccent]
    (0.2, 0.9) -- (5.3, 0.9)
    node[midway, above, lbl, text=DEAccent] {\tiny close()};

  % Shaded region showing data bypasses NameNode
  \fill[DEGreen!8] (0.15, 1.4) rectangle (11.1, 4.5);
  \node[font=\tiny\sffamily\itshape, text=DEGreen!60!black, rotate=90]
       at (11.3, 3.0) {data plane (bypasses NameNode)};

\end{tikzpicture}
\caption{HDFS read-path sequence diagram. Steps~1--2 are control-plane
         operations (RPC to the NameNode); steps~3--6 are data-plane
         operations (direct TCP to DataNodes). The shaded region highlights
         that all block data flows between client and DataNode, never
         through the NameNode. Checksum verification (step~5) runs at the
         client for every 512-byte chunk; a mismatch triggers a retry on
         a different replica.}
\label{fig:hdfs-read}
\end{figure}

The key performance characteristic of the read path is that data flows
directly between client and DataNode, bypassing the NameNode entirely.
This means that the NameNode's throughput is not a bottleneck for read
operations: it handles only the metadata lookup (typically measured in
microseconds), while data transfer runs at full network speed
(typically 1--10 Gbit/s per connection).

%--------------------------------------------------------------------
\subsection{The Write Path}
\label{subsec:write-path}
%--------------------------------------------------------------------

Writing to HDFS is more complex than reading, because data must be
reliably distributed to multiple DataNodes simultaneously while
ensuring consistency in the presence of potential DataNode failures.

\begin{enumerate}
  \item \textbf{Create.} The client calls \texttt{FileSystem.create()}.
        The client library contacts the NameNode to create a new
        file entry in the namespace. The NameNode verifies that the
        file does not already exist and that the client has write
        permission, then creates the file and grants the client a
        \term{write lease}.

  \item \textbf{Block allocation.} As the client accumulates data in
        a write buffer (default 64\KB\ per packet), it requests a new
        block from the NameNode. The NameNode selects DataNodes for
        the block---applying the rack-aware placement policy---and
        returns a list of DataNodes ordered by the pipeline sequence.

  \item \textbf{Pipeline construction.} The client establishes a
        TCP connection to the first DataNode in the pipeline; that
        DataNode connects to the second; the second connects to the
        third. This forms a \term{replication pipeline}: data flows
        through the pipeline rather than being sent from the client
        to each DataNode independently, reducing the client's outbound
        bandwidth requirement.

  \item \textbf{Packet streaming.} The client sends data in
        \term{packets} of 64\KB\ down the pipeline. As the first
        DataNode receives a packet, it writes it to its local disk and
        simultaneously forwards it to the second DataNode. The second
        DataNode writes and forwards to the third. Each DataNode sends
        an \term{acknowledgement} back up the pipeline once it has
        durably written the packet.

  \item \textbf{Block completion.} When all packets for a block have
        been acknowledged by all three DataNodes, the client notifies
        the NameNode that the block is complete. The NameNode records
        the block's DataNode locations in its metadata.

  \item \textbf{Close.} When the client calls \texttt{close()}, all
        remaining buffered packets are flushed through the pipeline,
        all blocks are completed, and the NameNode marks the file as
        complete, releasing the write lease.
\end{enumerate}

\begin{figure}[ht]
\centering
\begin{tikzpicture}[
    comp/.style={rectangle, rounded corners=3pt,
                 minimum width=2.0cm, minimum height=0.85cm,
                 font=\small\bfseries\sffamily, align=center},
    arr/.style={-{Stealth[length=5pt]}, thick},
    lbl/.style={font=\scriptsize\sffamily, align=center}
  ]
  \node[comp, fill=DELightAccent, draw=DEAccent!70] (client) at (0,   2.5) {Client};
  \node[comp, fill=DELightBlue,   draw=DEBlue!50]   (nn)     at (0,   0)   {NameNode};
  \node[comp, fill=DELightGreen,  draw=DEGreen!60]  (dn1)    at (4.5, 2.5) {DataNode 1\\{\scriptsize\normalfont(Rack A)}};
  \node[comp, fill=DELightGreen,  draw=DEGreen!60]  (dn2)    at (7.5, 2.5) {DataNode 2\\{\scriptsize\normalfont(Rack B)}};
  \node[comp, fill=DELightGreen,  draw=DEGreen!60]  (dn3)    at (10.5,2.5) {DataNode 3\\{\scriptsize\normalfont(Rack B)}};

  % Control plane: client <-> NameNode
  \draw[arr, DEBlue] (client.south) -- (nn.north)
        node[midway, left, lbl] {1. create()\\3. add block\\7. complete()};

  % Data pipeline
  \draw[arr, DEGreen!70!black, very thick]
        (client.east) -- (dn1.west)
        node[midway, above, lbl] {4. packets};
  \draw[arr, DEGreen!70!black, very thick]
        (dn1.east) -- (dn2.west)
        node[midway, above, lbl] {forward};
  \draw[arr, DEGreen!70!black, very thick]
        (dn2.east) -- (dn3.west)
        node[midway, above, lbl] {forward};

  % ACK pipeline (dashed, reverse direction)
  \draw[arr, DEAccent!80, dashed]
        (dn3.south) to[out=-90,in=-90,looseness=0.5] (dn2.south)
        node[midway, below, lbl] {ACK};
  \draw[arr, DEAccent!80, dashed]
        (dn2.south) to[out=-90,in=-90,looseness=0.6] (dn1.south)
        node[midway, below, lbl] {ACK};
  \draw[arr, DEAccent!80, dashed]
        (dn1.south) to[out=-90,in=-90,looseness=1.2] (client.south)
        node[midway, below, lbl] {ACK};

  % Labels
  \node[lbl, text=DEBlue]  at (0, -1.0) {Control plane (RPC)};
  \node[lbl, text=DEGreen!80!black] at (7.5, 3.5) {Data pipeline (TCP)};
\end{tikzpicture}
\caption{HDFS write path. The client communicates with the NameNode
         only for control messages; all block data flows through the
         three-DataNode replication pipeline. ACKs propagate back
         through the pipeline to confirm durable writes.}
\label{fig:hdfs-write}
\end{figure}

\begin{systemdesign}{Pipeline Failure Handling}
If a DataNode fails mid-pipeline during a write, HDFS recovers as
follows. The pipeline is closed. Packets that were acknowledged by
the failed DataNode but not yet confirmed by the client are marked for
re-send. A new block replica is allocated on a healthy DataNode, and a
new pipeline is constructed excluding the failed DataNode. The NameNode
is notified of the incomplete replica on the failed node and schedules
re-replication after the write completes. The client experiences a
brief pause but the write succeeds as long as at least one DataNode
survives. This mechanism allows HDFS to absorb DataNode failures
\emph{during} a write without requiring the client to restart the
entire write operation.
\end{systemdesign}

%%====================================================================
\section{HDFS High Availability}
\label{sec:hdfs-ha}
%%====================================================================

In the original HDFS design, the NameNode is a \term{single point of
failure} (SPOF). If the NameNode process crashes or the NameNode
machine becomes unavailable, the entire HDFS cluster is inaccessible:
no client can resolve file names to block locations, and no DataNode
can be instructed to re-replicate data. At the scale of Yahoo!'s
production Hadoop cluster---42,000 nodes in 2011---an unplanned
NameNode outage meant that 42,000 DataNodes became inaccessible
simultaneously. The resulting loss of analytical capacity during the
recovery period (which could take 30--60 minutes for a large
namespace) was a significant operational risk.

Hadoop 2.0 (released 2012) introduced \term{HDFS High Availability}
(HA), which eliminates the NameNode SPOF by maintaining two NameNode
instances simultaneously: one \term{Active NameNode} that handles all
client and DataNode traffic, and one \term{Standby NameNode} that
continuously mirrors the Active's metadata state and can take over
within seconds if the Active fails.

%--------------------------------------------------------------------
\subsection{The HA Architecture}
\label{subsec:ha-arch}
%--------------------------------------------------------------------

The HA architecture requires three key components beyond the standard
HDFS setup.

\textbf{Journal Nodes} are a quorum of lightweight processes (a minimum
of three, to tolerate one failure) that implement a shared,
fault-tolerant edit log. The Active NameNode writes every namespace
mutation to the Journal Node quorum as it executes; the Standby
NameNode reads from the Journal Node quorum continuously, replaying
edits to keep its in-memory namespace state synchronised with the
Active. Because the Journal Node quorum is shared, the Standby's
namespace state is always at most a few hundred milliseconds behind
the Active.

\textbf{ZooKeeper} provides distributed leader election. Each NameNode
runs a \term{ZooKeeper Failover Controller} (ZKFC) process that
monitors the NameNode's health and maintains a session with ZooKeeper.
If the Active NameNode becomes unhealthy---as detected by the ZKFC---
the Active's ZooKeeper session expires, and ZooKeeper triggers a
leader election. The Standby's ZKFC wins the election and promotes the
Standby to Active.

\textbf{Fencing} prevents the \term{split-brain} scenario---where both
NameNodes believe themselves to be Active and independently accept
writes, resulting in divergent metadata states. Before promoting the
Standby, the ZKFC executes a fencing operation that either gracefully
stops the former Active (via SSH) or, if the former Active is
unreachable, forcibly terminates its ability to communicate with
DataNodes using SSH-based remote kill commands or STONITH
(``Shoot The Other Node In The Head'') hardware mechanisms.

\begin{table}[ht]
\centering
\caption{NameNode HA component summary}
\label{tab:ha-components}
\renewcommand{\arraystretch}{1.4}
\begin{tabular}{L{2.8cm} L{7.8cm}}
\toprule
\textbf{Component} & \textbf{Role} \\
\midrule
Active NameNode    & Handles all client and DataNode traffic; writes edits to JN quorum \\
Standby NameNode   & Continuously replays JN edits; ready to take over within seconds \\
Journal Nodes (3+) & Shared, quorum-based edit log; tolerates minority failures \\
ZooKeeper (3+)     & Leader election and health monitoring for the two NameNodes \\
ZKFC               & Per-NameNode daemon that manages ZooKeeper session and triggers failover \\
\bottomrule
\end{tabular}
\end{table}

\begin{figure}[ht]
\centering
\begin{tikzpicture}[
    ann/.style={rectangle, draw=DEBlue, fill=DELightBlue!50, rounded corners=4pt,
                minimum width=3.0cm, minimum height=1.2cm, align=center,
                font=\small\sffamily},
    snn/.style={rectangle, draw=DEBlue!60, fill=DELightBlue!25, rounded corners=4pt,
                minimum width=3.0cm, minimum height=1.2cm, align=center,
                font=\small\sffamily},
    jn/.style={rectangle, draw=DEGreen!70!black, fill=DELightGreen, rounded corners=3pt,
               minimum width=1.6cm, minimum height=0.85cm, align=center,
               font=\scriptsize\sffamily},
    zk/.style={rectangle, draw=DEAccent!80, fill=DELightAccent, rounded corners=3pt,
               minimum width=1.5cm, minimum height=0.85cm, align=center,
               font=\scriptsize\sffamily},
    zkfc/.style={rectangle, draw=DEGray!60, fill=DEGray!12, rounded corners=2pt,
                 minimum width=1.4cm, minimum height=0.55cm, align=center,
                 font=\tiny\sffamily},
    dn/.style={rectangle, draw=DEGray!60, fill=DEGray!8, rounded corners=3pt,
               minimum width=1.6cm, minimum height=0.7cm, align=center,
               font=\scriptsize\sffamily},
    arr/.style={-{Stealth[length=4pt,width=4pt]}, thick},
    darr/.style={{Stealth[length=4pt,width=4pt]}-{Stealth[length=4pt,width=4pt]}, thick},
    lbl/.style={font=\scriptsize\sffamily, align=center}
  ]

  %--- Active NameNode (left) ---
  \node[ann] (ann) at (1.8, 7.5) {Active\\NameNode};
  \node[zkfc] (zkfc1) at (1.8, 6.6) {ZKFC};

  %--- Standby NameNode (right) ---
  \node[snn] (snn) at (9.2, 7.5) {Standby\\NameNode};
  \node[zkfc] (zkfc2) at (9.2, 6.6) {ZKFC};

  %--- Journal Nodes (middle row) ---
  \node[jn] (jn1) at (3.9, 7.5) {Journal\\Node 1};
  \node[jn] (jn2) at (5.5, 7.5) {Journal\\Node 2};
  \node[jn] (jn3) at (7.1, 7.5) {Journal\\Node 3};
  \node[font=\tiny\sffamily\color{DEGreen!60!black}] at (5.5, 8.4)
       {Journal Node Quorum (shared EditLog)};
  \draw[draw=DEGreen!40, dashed, rounded corners=3pt]
       (3.2, 7.1) rectangle (7.8, 8.2);

  %--- ZooKeeper cluster (middle, lower) ---
  \node[zk] (zk1) at (3.9, 5.5) {ZooKeeper\\Node 1};
  \node[zk] (zk2) at (5.5, 5.5) {ZooKeeper\\Node 2};
  \node[zk] (zk3) at (7.1, 5.5) {ZooKeeper\\Node 3};
  \node[font=\tiny\sffamily\color{DEAccent!80}] at (5.5, 6.5)
       {ZooKeeper Cluster (leader election)};
  \draw[draw=DEAccent!35, dashed, rounded corners=3pt]
       (3.1, 5.1) rectangle (7.9, 6.3);

  %--- DataNodes (bottom row, report to BOTH NNs) ---
  \node[dn] (dn1) at (2.2, 3.6) {DataNode 1};
  \node[dn] (dn2) at (4.2, 3.6) {DataNode 2};
  \node[dn] (dn3) at (6.2, 3.6) {DataNode 3};
  \node[dn] (dn4) at (8.2, 3.6) {DataNode $N$};
  \node[font=\tiny\sffamily\color{DEGray!60}] at (5.2, 3.1)
       {DataNodes (heartbeat + block report to both NameNodes)};

  %--- Active NN writes to Journal Nodes ---
  \draw[arr, DEGreen!70!black, thick]
    (ann.east) -- (jn1.west)
    node[midway, above, lbl, text=DEGreen!70!black] {\tiny edits};

  %--- Standby NN reads from Journal Nodes ---
  \draw[arr, DEGreen!70!black, thick]
    (jn3.east) -- (snn.west)
    node[midway, above, lbl, text=DEGreen!70!black] {\tiny replay};

  %--- ZKFC <-> ZooKeeper (both NNs) ---
  \draw[darr, DEAccent!70]
    (zkfc1.south) to[out=-60, in=135] (zk1.north west)
    node[pos=0.4, left, lbl, text=DEAccent!80] {\tiny session};
  \draw[darr, DEAccent!70]
    (zkfc2.south) to[out=220, in=45] (zk3.north east)
    node[pos=0.4, right, lbl, text=DEAccent!80] {\tiny session};

  %--- DataNodes -> Active NN (heartbeat) ---
  \draw[darr, DEGray!60]
    (dn1.north) to[out=100, in=-130] (ann.south)
    node[pos=0.3, left, lbl, text=DEGray!60] {\tiny HB};
  %--- DataNodes -> Standby NN (heartbeat) ---
  \draw[darr, DEGray!60]
    (dn4.north) to[out=80, in=-50] (snn.south)
    node[pos=0.3, right, lbl, text=DEGray!60] {\tiny HB};
  \draw[darr, DEGray!55]
    (dn2.north) to[out=120, in=-100] (ann.south);
  \draw[darr, DEGray!55]
    (dn3.north) to[out=60, in=-80] (snn.south);

  %--- Fencing arrow (ZKFC1 fences Active before Standby promotes) ---
  \draw[arr, DERed!70, very thick, dashed]
    (zkfc2.north) to[out=90, in=-90] (zkfc1.south)
    node[pos=0.5, left, lbl, text=DERed!70]
         {\tiny fencing\\[-2pt]\tiny (on failover)};

  %--- Legend ---
  \node[font=\tiny\sffamily, align=left, text=DEGray] at (11.2, 6.8) {%
    \textbf{Legend}\\[2pt]
    \textcolor{DEGreen!70!black}{$\rightarrow$} EditLog flow\\
    \textcolor{DEAccent!80}{$\leftrightarrow$} ZooKeeper session\\
    \textcolor{DEGray!60}{$\leftrightarrow$} Heartbeat\\
    \textcolor{DERed!70}{- -\!-\!>} Fencing signal};

\end{tikzpicture}
\caption{HDFS High Availability architecture. The Active NameNode writes every
         namespace mutation to the Journal Node quorum (shared EditLog); the
         Standby NameNode continuously replays those edits to maintain a
         synchronised copy of the namespace. Both NameNodes run a
         ZooKeeper Failover Controller (ZKFC) that monitors health and
         maintains a ZooKeeper session. On Active failure, the ZKFC fences
         the former Active (preventing split-brain), then promotes the Standby
         to Active in seconds. DataNodes send heartbeats and block reports to
         \emph{both} NameNodes, so the Standby has a live block map
         ready for immediate service.}
\label{fig:hdfs-ha}
\end{figure}

%%====================================================================
\section{HDFS Design Characteristics and Limitations}
\label{sec:hdfs-limits}
%%====================================================================

HDFS is a purpose-built system optimised for specific workloads. Like
GFS, its strengths and weaknesses are two sides of the same coin: the
design decisions that make it excellent for batch processing of large
files make it unsuitable for other access patterns.

%--------------------------------------------------------------------
\subsection{What HDFS Is Optimised For}
\label{subsec:hdfs-strengths}
%--------------------------------------------------------------------

HDFS excels at \textbf{high-throughput, sequential access to large
files}. A well-tuned HDFS cluster can deliver aggregate read throughput
of several GB/s per node: with 100 DataNodes each delivering 200 MB/s
from local disk, the cluster sustains 20 GB/s of aggregate read
bandwidth. For workloads that need to scan entire datasets---MapReduce
jobs, Hive queries that process all records in a table, Spark jobs that
re-partition a large dataset---this throughput is the critical metric,
and HDFS is optimised for it.

HDFS also provides \textbf{fault-tolerant write durability}. A write
that returns a success acknowledgement to the client is guaranteed to
have been committed to at least $r$ DataNodes (where $r$ is the
replication factor). Hardware failure of any subset of up to $r-1$
DataNodes simultaneously cannot cause data loss for any successfully
written block.

%--------------------------------------------------------------------
\subsection{The Small File Problem}
\label{subsec:small-file}
%--------------------------------------------------------------------

HDFS's single most significant operational limitation is its poor
handling of large numbers of small files. Each file, regardless of
size, consumes approximately 150 bytes of NameNode heap space for its
metadata. A single large file split into ten 128\MB\ blocks consumes
approximately 1.65\KB\ of NameNode memory. Ten million small files---
each taking one block---consume 1.5\GB\ of NameNode heap, regardless
of the files' total data volume. At 100 million small files, the
NameNode requires 15\GB\ of heap and metadata operations begin to slow
noticeably.

This is a critical practical constraint. Web server log files,
sensor telemetry archives, and social media ingestion pipelines
naturally produce large numbers of small files if not carefully managed.
The standard engineering response is to use \term{SequenceFiles}---
container files that aggregate many small records into a single large
HDFS file, providing a sequential index for record-level access---
or to pre-aggregate small files before ingestion using tools such as
HDFS \texttt{archive} (HAR files) or Apache CombineFileInputFormat
in the MapReduce layer.

%--------------------------------------------------------------------
\subsection{No In-Place Random Writes}
\label{subsec:no-random-write}
%--------------------------------------------------------------------

HDFS does not support random writes to existing files. Once a block
has been written and acknowledged, its content cannot be modified
in place. HDFS supports only sequential writes to new files and
atomic appends to the end of existing files (the append semantics
introduced in HDFS 0.21 and Hadoop 2.x). This is a deliberate
architectural choice that greatly simplifies the consistency model and
the replication protocol: if blocks are immutable after writing, the
system never needs to coordinate writes across multiple replicas.

For workloads that require in-place updates---a transactional database,
an incremental index---HDFS is the wrong tool. Apache HBase (examined
in Chapter~\ref{ch:mapreduce}) provides random read/write access on
top of HDFS by using an LSM-tree (Log-Structured Merge-tree) structure
that converts random writes into sequential HDFS writes.

%--------------------------------------------------------------------
\subsection{No Low-Latency Random Access}
\label{subsec:no-random-read}
%--------------------------------------------------------------------

HDFS is tuned for throughput, not latency. Opening an HDFS file
requires a round trip to the NameNode; reading a specific block
requires a TCP connection setup to the appropriate DataNode. Total
latency for a random read is typically in the range of 50--200 ms---
acceptable for batch processing but wholly unacceptable for
interactive queries (which require sub-millisecond storage access)
or for OLTP workloads. This is why HDFS is used as the storage
substrate for batch and near-batch systems (MapReduce, Spark, Hive),
not as a replacement for relational databases or key-value stores.

\begin{table}[ht]
\centering
\caption{HDFS strengths, limitations, and engineering mitigations}
\label{tab:hdfs-limits}
\renewcommand{\arraystretch}{1.45}
\begin{tabular}{L{3.2cm} L{3.4cm} L{4.0cm}}
\toprule
\textbf{Characteristic} & \textbf{Impact} & \textbf{Mitigation / alternative} \\
\midrule
Optimised for sequential reads & Excellent batch throughput & Use for MapReduce, Spark, Hive \\
Large block size (128 MB)      & Poor for small files      & SequenceFiles, HAR archives, CombineFileInputFormat \\
No random writes               & Cannot update in place    & HBase (LSM-tree over HDFS) \\
High read latency (50--200 ms) & No interactive queries    & Presto, Spark SQL with in-memory caching \\
NameNode SPOF (pre-HA)         & Cluster-wide outage       & HDFS HA with ZooKeeper \\
Namespace scalability          & 100M file ceiling         & HDFS Federation (multiple NameNodes) \\
\bottomrule
\end{tabular}
\end{table}

%%====================================================================
\section{Beyond HDFS: Cloud Object Storage}
\label{sec:object-storage}
%%====================================================================

Starting around 2012 and accelerating through 2015--2018, the data
engineering industry began a significant migration away from
on-premises HDFS clusters toward cloud-based \term{object storage}:
Amazon S3, Google Cloud Storage (GCS), and Azure Data Lake Storage
(ADLS). By 2024, cloud object storage has largely supplanted HDFS as
the primary storage layer for new data engineering deployments, though
HDFS remains in use at organisations that operate their own data
centres.

Cloud object storage differs from HDFS in several important ways.
Objects in S3 or GCS are identified by a key (essentially a path
string) and stored as a flat namespace with no directory hierarchy
(though object key prefixes mimic directory structure). Unlike HDFS
blocks, objects can be up to 5 TB each. Object storage is practically
unlimited in capacity and scales automatically with usage---there is
no NameNode to run out of heap, no replication factor to configure
manually, and no DataNodes to provision and maintain.

From a pricing and operations standpoint, cloud object storage is
dramatically simpler. Amazon S3 charges approximately \$0.023/GB/month
for standard storage; there are no servers to operate, no replication
factor to configure, and no NameNode HA to maintain. A data engineering
team using S3 eliminates the entire operational burden of the HDFS
cluster in exchange for a predictable per-GB-per-month cost.

The principal trade-off is \textbf{decoupling storage from compute}.
In a Hadoop cluster, DataNodes run both storage and compute (MapReduce
or Spark tasks), exploiting data locality. Cloud object storage is
physically separate from the compute cluster; every read requires a
network hop. However, cloud data centre networks now deliver bandwidth
of 25--100 Gbit/s between storage and compute instances, which is
often higher than the inter-DataNode bandwidth in on-premises
Hadoop clusters. Modern frameworks such as Apache Spark have been
optimised to run effectively against object storage, and the loss of
data locality has proven to be a minor performance penalty in most
workloads compared to the operational simplicity of S3.

\begin{keyinsight}{The Shift to Storage-Compute Separation}
The move from HDFS to cloud object storage represents a fundamental
architectural shift: from \emph{tightly coupled} storage-compute
(where data lives next to the CPU that processes it) to \emph{loosely
coupled} (where storage and compute scale independently). In a Hadoop
cluster, you cannot add compute capacity without also adding storage,
and vice versa. With S3 + cloud compute, you can scale a Spark cluster
from 10 to 1,000 nodes and back to 10 without moving a single byte
of data. This elasticity has proven to be the dominant engineering
advantage, and it is the primary reason cloud object storage has
displaced on-premises HDFS for new deployments.
\end{keyinsight}

%%====================================================================
\section{Case Study: Yahoo!'s Hadoop Cluster (2006--2011)}
\label{sec:yahoo-case}
%%====================================================================

\begin{casestudy}{Yahoo! --- The First Industrial-Scale Open-Source Distributed File System}

\textbf{Background.}
Yahoo!\ was the first organisation outside Google to operate a
production Hadoop cluster at true industrial scale. The company had
several workloads that required petabyte-scale batch processing: web
ranking (recomputing the search index nightly from the full crawl),
user behaviour analytics (processing billions of page views per day),
and spam detection (scanning email metadata across hundreds of millions
of user accounts). After evaluating alternatives, Yahoo!\ bet its
entire search and analytics infrastructure on the open-source Hadoop
ecosystem starting in 2006.

\smallskip
\textbf{Scale and growth.}
Yahoo!'s Hadoop cluster grew from approximately 1,000 nodes in 2006
to a peak of \textbf{42,000 nodes} in 2011---a 42-fold scale-up in
five years. At its peak, the cluster stored approximately \textbf{100
petabytes} of data in HDFS across two major data centres. The web
index computation alone ran tens of thousands of MapReduce jobs per
day. Yahoo!'s infrastructure team reported that the cluster
experienced approximately \textbf{2--3 DataNode failures per day} on
average---exactly the failure rate predicted by the GFS paper for
commodity hardware at this scale.

\smallskip
\textbf{Operational challenges.}
Managing a 42,000-node HDFS cluster required solving engineering
problems that the original Hadoop design had not anticipated.

The \emph{NameNode bottleneck} emerged as the cluster grew. By 2009,
Yahoo!'s HDFS namespace contained over 50 billion blocks, consuming
the full 64\GB\ of available NameNode heap. Startup time after a
planned NameNode maintenance window exceeded 90 minutes---during
which the entire cluster was inaccessible. Yahoo!'s engineering team
developed several mitigations, including a distributed NameNode
extension (\emph{HDFS Federation}) that partitioned the namespace
across multiple NameNodes, reducing the per-NameNode block count and
memory pressure.

The \emph{small file problem} was acute in production. Yahoo!'s log
ingestion pipelines initially wrote one file per server per minute,
producing hundreds of millions of small files per day. The engineering
team developed automated compaction pipelines that aggregated hourly
log files into daily SequenceFiles, reducing the namespace size by a
factor of roughly 1,000 while preserving query access.

\smallskip
\textbf{Contributions to open source.}
Yahoo!\ was by far the largest contributor to the Apache Hadoop
project during this period. Key features that originated from Yahoo!'s
production experience include: HDFS Append (enabling streaming log
ingestion), Hadoop Security (Kerberos-based authentication, developed
after Yahoo!\ experienced data governance incidents with an unsecured
cluster), HDFS High Availability (eliminating the NameNode SPOF after
a NameNode failure caused a 45-minute cluster-wide outage in 2010),
and YARN (the resource manager that replaced the original
MapReduce-only scheduler, enabling non-MapReduce workloads such as
Spark and Tez to share the cluster).

\smallskip
\textbf{Lessons for data engineering.}
Yahoo!'s experience establishes several principles that remain
relevant today:
\begin{itemize}
  \item HDFS is exceptionally robust for large-file, sequential-access
        workloads at any scale---but it requires careful operational
        management, particularly of the NameNode's memory and the
        accumulation of small files.
  \item Production distributed systems require active contributions
        back to their upstream open-source projects: Yahoo!\ could not
        have operated at 42,000 nodes without the HDFS improvements it
        developed and contributed. This model of ``upstream dependency,
        active contribution'' has been adopted by Facebook, LinkedIn,
        Twitter (X), and most major technology companies.
  \item Failure planning is not optional. Yahoo!'s decision to invest
        in HDFS HA before a major outage occurred---rather than after---
        is the textbook example of proactive operational engineering.
\end{itemize}
\end{casestudy}

%%====================================================================
\section*{Chapter Summary}
\label{sec:ch3-summary}
%%====================================================================

\begin{itemize}
  \item Conventional file systems (NFS, SAN, POSIX) fail at petabyte
        scale for three reasons: single-machine capacity limits,
        insufficient sequential throughput, and an inability to tolerate
        the continuous hardware failures that occur at cluster scale.

  \item \textbf{GFS} (2003) solved these problems by: using large 64\MB\
        chunks, maintaining all metadata in a single master's RAM,
        separating control-plane (master) from data-plane (chunkservers),
        supporting append-only semantics, and relying on replication
        rather than RAID for fault tolerance.

  \item \textbf{HDFS} is the open-source GFS implementation. Its key
        components are the \textbf{NameNode} (namespace and block mapping
        in RAM, EditLog + FsImage persistence) and \textbf{DataNodes}
        (block storage, heartbeat, block report, checksum verification).

  \item \textbf{Rack-aware replication} (factor 3, default) places one
        replica locally, one on a different rack, one on a second node in
        the same rack as replica 2. A single rack failure cannot cause
        data loss.

  \item The \textbf{write pipeline} flows client $\to$ DN1 $\to$ DN2
        $\to$ DN3 with acknowledgements propagating in reverse. The
        NameNode handles only control messages; data bypasses it entirely.

  \item \textbf{HDFS HA} uses Active/Standby NameNodes, shared Journal
        Nodes for the edit log, ZooKeeper for leader election, and
        fencing to prevent split-brain.

  \item HDFS limitations include the \textbf{small file problem}
        (NameNode heap exhaustion), \textbf{no random writes}, and
        \textbf{high read latency}. Mitigations include SequenceFiles,
        HBase, and Presto.

  \item \textbf{Cloud object storage} (S3, GCS, ADLS) has largely
        replaced HDFS for new deployments by offering unlimited capacity,
        no operational burden, and elastic compute-storage separation---at
        the cost of data locality.
\end{itemize}

%%====================================================================
\section*{Exercises}
\label{sec:ch3-exercises}
%%====================================================================

\exercisesection{Short Questions}

\sq State three reasons why a conventional NFS-based file system cannot
store and process a 500-petabyte dataset.

\sq What is a GFS chunk? What is the default size and why was this size
chosen instead of the 4-\KB\ blocks used by typical file systems?

\sq Explain the GFS design decision to keep all metadata in the master's
RAM. What is the practical ceiling on the number of chunks this imposes,
and at what dataset size is that ceiling reached for 64-\MB\ chunks?

\sq What is the HDFS NameNode? List its five responsibilities and
explain what would happen to a running HDFS cluster if the NameNode
process were killed.

\sq What is the difference between the HDFS Secondary NameNode and a
hot-standby NameNode? What problem does the Secondary NameNode actually
solve?

\sq Describe the rack-aware replication policy for replication factor 3.
If a rack fails, how many replicas of each block remain? Does this
satisfy the fault-tolerance goal?

\sq Trace the HDFS write path step by step from \texttt{FileSystem.create()}
to the final acknowledgement. At which step is the data actually sent
to DataNodes, and how many network hops does it take before the client
receives a write acknowledgement?

\sq What is a DataNode heartbeat and what is a block report? How does the
NameNode use each to manage the cluster?

\sq Explain the HDFS small file problem. If a 100-node cluster holds
200 million files, each occupying one 128-\MB\ block, estimate the
minimum NameNode heap required (assume 300 bytes per file--block pair).

\sq Why does HDFS not support random writes? What are the consistency
advantages of the immutable-block design?

\sq What is HDFS High Availability? List the three components beyond the
standard HDFS setup that HA requires and explain the role of each.

\sq What is a split-brain condition in an HA system? Why is fencing
necessary, and what are two mechanisms by which HDFS HA can fence the
former Active NameNode?

\sq Compare HDFS and Amazon S3 on four dimensions: (a) data locality,
(b) operational burden, (c) capacity ceiling, (d) cost model.

\sq A data engineering team ingests one million web server log files per
day, each 200\KB\ in size, directly into HDFS. After one year, the
NameNode heap is full. Identify the root cause and propose two
engineering solutions.

\bigskip
\exercisesection{Analytical Problems}

\ap \textbf{Replication and Storage Overhead Analysis.}
A genomics research institute stores 50 petabytes of raw sequencing
data in HDFS with the default replication factor of 3 and a block
size of 128\MB.

\begin{enumerate}[label=(\alph*)]
  \item Calculate the total HDFS storage capacity required (including
        replicas) in petabytes.
  \item Calculate the number of 128-\MB\ blocks required to store
        50\PB\ of data.
  \item If each DataNode has 200\TB\ of usable disk capacity, how many
        DataNodes are required to store the replicated dataset?
  \item Estimate the NameNode heap required to hold metadata for all
        blocks (assume 150 bytes per block). Is this within the 128\GB\
        of RAM on a modern server?
  \item If the institute reduces the replication factor to 2, what is
        the storage saving in petabytes? What is the trade-off?
\end{enumerate}

\ap \textbf{HDFS Write Throughput Analysis.}
A streaming ingestion pipeline writes data to HDFS at a sustained rate
of 2\GB/s. Each DataNode has a 10 Gbit/s network interface (effective
throughput 1.25\GB/s). The replication factor is 3.

\begin{enumerate}[label=(\alph*)]
  \item In the write pipeline (client $\to$ DN1 $\to$ DN2 $\to$ DN3),
        how much network bandwidth does each hop consume at 2\GB/s
        ingestion rate?
  \item Is a single DataNode's network interface the bottleneck? If the
        pipeline is the only network traffic, at what ingestion rate
        would the first DataNode's network interface saturate?
  \item How many simultaneous parallel write pipelines (each ingesting
        a different block) can a 10 Gbit/s interface support at 2\GB/s
        per pipeline? Why is it beneficial to parallelise writes?
  \item At 2\GB/s ingestion rate and 128\MB\ blocks, how many blocks
        per second does the system create? At what rate does the
        NameNode receive \texttt{addBlock} RPC calls?
\end{enumerate}

\ap \textbf{NameNode HA Failover Timing.}
An HDFS HA cluster has the following configuration:
\begin{itemize}
  \item ZooKeeper session timeout: 10 seconds
  \item ZKFC health check interval: 5 seconds
  \item Standby NameNode edit log replay lag: 200 milliseconds maximum
  \item SSH fencing timeout: 30 seconds
\end{itemize}

\begin{enumerate}[label=(\alph*)]
  \item What is the minimum time from Active NameNode failure to the
        Standby becoming Active (assuming fencing via SSH succeeds)?
  \item What is the maximum time in the worst case (ZooKeeper session
        timeout has just reset, fencing takes the full 30 seconds)?
  \item During the failover window, what happens to HDFS clients that
        are actively reading? What happens to clients that are in the
        middle of a write?
  \item Explain why fencing must be completed before the Standby is
        promoted. Describe the consequences of a successful promotion
        without fencing, given that the former Active's write lease is
        still valid.
\end{enumerate}

\ap \textbf{Cloud vs.\ On-Premises Storage Economics.}
A media company stores 2\PB\ of video content in on-premises HDFS
with replication factor 3 (requiring 6\PB\ raw disk capacity). The
data grows at 500\TB/year. DataNode servers cost \$8,000 each with
72\TB\ usable disk per server; server lifetime is 4 years (amortised
\$2,000/server/year). Operations staff costs \$150,000/year to manage
the cluster. Amazon S3 charges \$0.023/GB/month for standard storage.

\begin{enumerate}[label=(\alph*)]
  \item How many DataNode servers are needed today? What is the
        annualised hardware cost (amortised over 4 years)?
  \item What is the total annual on-premises cost (hardware + operations)?
  \item What is the annual S3 cost for 2\PB\ of raw data (no
        replication factor needed---S3 manages redundancy internally)?
  \item After 5 years of 500\TB/year growth, compare the cumulative
        5-year costs of both approaches (assume S3 pricing remains
        constant; on-premises requires a hardware refresh at year 4).
  \item What non-monetary factors should also influence this decision?
\end{enumerate}

\bigskip
\exercisesection{Design Problems}

\dpr \textbf{HDFS Cluster Design for a Video Streaming Platform.}
A video streaming company archives 8\PB\ of source video files for
transcoding and long-term retention. The files range from 500\MB\ to
40\GB\ each (typical feature film at broadcast quality). The company
also processes 2 TB/day of viewer analytics events (one JSON record
per play event, approximately 512 bytes each, producing $\sim$4
billion records per day). The system must sustain 10\GB/s sustained
write throughput during peak ingestion.

Design the HDFS cluster and ingestion pipeline. Your design must specify:
\begin{enumerate}[label=(\alph*)]
  \item The appropriate replication factor for video files vs.\ analytics
        events, and the justification for each choice.
  \item The HDFS block size configuration for video files and for
        analytics events (can differ per dataset).
  \item The DataNode hardware specification (disk capacity, network
        interface speed) and the number of DataNodes required.
  \item How the 4-billion-record-per-day analytics stream should be
        ingested without creating the small file problem (propose a
        specific ingestion architecture using tools from this or
        previous chapters).
  \item The HA configuration for the NameNode, including the minimum
        number of Journal Nodes and ZooKeeper nodes required for
        quorum.
  \item The monitoring metrics you would track in production and the
        alerting thresholds that would trigger an on-call escalation.
\end{enumerate}

\dpr \textbf{HDFS Migration to Cloud Object Storage.}
A financial services company operates a 500-node on-premises HDFS
cluster storing 1.2\PB\ of trade data, customer records, and
regulatory filings. The company is evaluating a migration to Amazon
S3. The data has strict regulatory requirements: all data must be
encrypted at rest and in transit, all access must be auditable (for
SEC and FINRA compliance), and data must be retained for 7 years.
Certain datasets must remain on-premises (customer PII, subject to
contractual data residency obligations).

Design the migration strategy. Your design must address:
\begin{enumerate}[label=(\alph*)]
  \item A data classification scheme that determines which datasets
        can migrate to S3 and which must remain on-premises.
  \item The migration mechanism: how to transfer 1.2\PB\ to S3 with
        minimal disruption to running production pipelines.
  \item How encryption at rest and in transit is enforced for S3-stored
        data.
  \item How access auditing (``who accessed which object, when'') is
        implemented in S3.
  \item The hybrid architecture that allows on-premises Spark jobs to
        read data from both the remaining HDFS cluster and S3
        transparently.
  \item The rollback plan if the S3 migration reveals unacceptable
        latency or data access pattern incompatibilities.
\end{enumerate}

\bigskip
\exercisesection{Programming Exercises}

\pe \textbf{HDFS Block Layout Simulator (Python).}
Simulate the block layout for a set of files in HDFS: compute how many
blocks each file occupies, how much space is wasted in the last block,
and the total NameNode metadata overhead.

\begin{lstlisting}[style=pyspark, caption={PE3.1: HDFS block layout simulator}]
BLOCK_SIZE_MB = 128
BYTES_PER_MB = 1024 * 1024
METADATA_BYTES_PER_BLOCK = 150  # NameNode heap per block

files = [
    {"name": "web_crawl_2024_01.seq",       "size_gb": 312.0},
    {"name": "user_events_2024_01_01.json",  "size_mb": 0.8},
    {"name": "product_images.tar",           "size_gb": 1450.0},
    {"name": "transaction_log_hourly.orc",   "size_gb": 48.5},
    {"name": "genome_sample_001.fastq",      "size_gb": 220.0},
]

def analyse_file(name, size_bytes, block_size=BLOCK_SIZE_MB * BYTES_PER_MB,
                 replication=3):
    """
    Return a dict with:
      num_blocks       : number of HDFS blocks (ceil division)
      last_block_bytes : size of the final (partial) block
      wasted_bytes     : padding to fill the last block to block_size
                         (HDFS does NOT pad, so wasted_bytes = 0 always;
                          but report the theoretical waste if it did)
      metadata_bytes   : NameNode heap for this file's blocks
      replicated_bytes : total physical disk used (data * replication)
    """
    import math
    num_blocks = math.ceil(size_bytes / block_size)
    last_block_bytes = size_bytes % block_size or block_size
    # HDFS does not pad the last block
    wasted_bytes = 0
    metadata_bytes = num_blocks * METADATA_BYTES_PER_BLOCK
    replicated_bytes = size_bytes * replication
    return {
        "name": name,
        "size_bytes": size_bytes,
        "num_blocks": num_blocks,
        "last_block_bytes": last_block_bytes,
        "wasted_bytes": wasted_bytes,
        "metadata_bytes": metadata_bytes,
        "replicated_bytes": replicated_bytes,
    }

def summarise(files_spec):
    results = []
    for f in files_spec:
        if "size_gb" in f:
            size = int(f["size_gb"] * 1024 * BYTES_PER_MB)
        else:
            size = int(f["size_mb"] * BYTES_PER_MB)
        results.append(analyse_file(f["name"], size))

    total_blocks = sum(r["num_blocks"] for r in results)
    total_nn_mb  = sum(r["metadata_bytes"] for r in results) / BYTES_PER_MB
    total_rep_tb = sum(r["replicated_bytes"] for r in results) / (1024**4)

    print(f"{'File':<38} {'Blocks':>7} {'Last block':>12} {'NN heap':>10}")
    print("-" * 72)
    for r in results:
        lb_mb = r["last_block_bytes"] / BYTES_PER_MB
        nn_kb = r["metadata_bytes"] / 1024
        print(f"{r['name']:<38} {r['num_blocks']:>7,} "
              f"{lb_mb:>10.1f} MB {nn_kb:>8.1f} KB")
    print("-" * 72)
    print(f"{'TOTAL':<38} {total_blocks:>7,} "
          f"{'':>12} {total_nn_mb:>8.1f} MB")
    print(f"\nTotal replicated storage (3x): {total_rep_tb:.2f} TB")

summarise(files)
\end{lstlisting}

\pe \textbf{Rack-Aware Replica Placement Simulator (Python).}
Simulate the HDFS rack-aware replica placement policy for a set of
blocks being written to a cluster of DataNodes distributed across racks.

\begin{lstlisting}[style=pyspark, caption={PE3.2: Rack-aware placement simulator}]
import random

# Cluster topology: rack -> list of DataNode IDs
cluster = {
    "rack-A": ["dn01", "dn02", "dn03", "dn04", "dn05"],
    "rack-B": ["dn06", "dn07", "dn08", "dn09", "dn10"],
    "rack-C": ["dn11", "dn12", "dn13", "dn14", "dn15"],
}

def get_rack(dn_id, cluster):
    for rack, nodes in cluster.items():
        if dn_id in nodes:
            return rack
    return None

def place_replicas(writer_dn, cluster, replication=3):
    """
    Apply HDFS rack-aware placement for replication factor 3:
      R1: same node as writer (or writer's rack if writer not in cluster)
      R2: a node on a DIFFERENT rack from R1
      R3: a DIFFERENT node on the SAME rack as R2

    Returns list of (dn_id, rack) tuples, one per replica.
    """
    writer_rack = get_rack(writer_dn, cluster)
    if writer_rack is None:
        # Writer is not a DataNode; pick any node for R1
        writer_rack = random.choice(list(cluster.keys()))
        r1 = random.choice(cluster[writer_rack])
    else:
        r1 = writer_dn

    # R2: different rack
    other_racks = [r for r in cluster if r != writer_rack]
    rack2 = random.choice(other_racks)
    r2 = random.choice(cluster[rack2])

    # R3: different node on same rack as R2
    rack2_others = [n for n in cluster[rack2] if n != r2]
    r3 = random.choice(rack2_others)

    return [(r1, writer_rack), (r2, rack2), (r3, rack2)]

def simulate_writes(num_blocks, writer_dn, cluster):
    placements = []
    for i in range(num_blocks):
        replicas = place_replicas(writer_dn, cluster)
        placements.append(replicas)

    # Count blocks per DataNode
    from collections import Counter
    node_counts = Counter()
    for rep_list in placements:
        for dn, rack in rep_list:
            node_counts[dn] += 1

    print(f"Block placement for {num_blocks} blocks "
          f"(writer: {writer_dn}, rack: {get_rack(writer_dn, cluster)})\n")
    print(f"{'Block':>6}  R1              R2              R3")
    print("-" * 55)
    for i, rep in enumerate(placements[:10]):
        cols = [f"{dn} ({rack})" for dn, rack in rep]
        print(f"{i+1:>6}  {cols[0]:<16}{cols[1]:<16}{cols[2]}")
    if num_blocks > 10:
        print(f"  ... ({num_blocks - 10} more blocks)")

    print(f"\nDataNode block distribution (top 5):")
    for dn, count in node_counts.most_common(5):
        rack = get_rack(dn, cluster)
        print(f"  {dn} ({rack}): {count} blocks")

simulate_writes(num_blocks=50, writer_dn="dn03", cluster=cluster)
\end{lstlisting}

\pe \textbf{Failure Simulation and Re-replication Scheduler (Python).}
Simulate the NameNode's re-replication logic when DataNodes fail:
identify under-replicated blocks and schedule re-replication from
surviving replicas.

\begin{lstlisting}[style=pyspark, caption={PE3.3: Re-replication scheduler}]
# Initial block map: block_id -> list of DataNode IDs holding a replica
block_map = {
    "blk_001": ["dn01", "dn04", "dn07"],
    "blk_002": ["dn02", "dn05", "dn09"],
    "blk_003": ["dn01", "dn06", "dn11"],
    "blk_004": ["dn03", "dn07", "dn12"],
    "blk_005": ["dn01", "dn02", "dn13"],
    "blk_006": ["dn04", "dn08", "dn14"],
    "blk_007": ["dn02", "dn06", "dn10"],
}

REPLICATION_FACTOR = 3
HEALTHY_NODES = {"dn01", "dn02", "dn03", "dn04",
                 "dn05", "dn06", "dn07", "dn08",
                 "dn09", "dn10", "dn11", "dn12", "dn13", "dn14"}

def simulate_failure(failed_nodes, block_map, healthy_nodes,
                     replication=REPLICATION_FACTOR):
    """
    Remove failed_nodes from block_map and healthy_nodes.
    Return a list of re-replication tasks:
      (block_id, source_dn, priority)
    where priority = 'CRITICAL' (1 surviving replica),
                     'HIGH'     (2 surviving replicas),
                     'NORMAL'   (0 deficit, should not appear)
    Tasks are sorted highest priority first.
    """
    live_nodes = healthy_nodes - set(failed_nodes)
    tasks = []

    for blk_id, replicas in block_map.items():
        surviving = [r for r in replicas if r in live_nodes]
        deficit = replication - len(surviving)
        if deficit <= 0:
            continue  # no re-replication needed
        priority = "CRITICAL" if len(surviving) == 1 else "HIGH"
        # Source: first surviving replica
        source = surviving[0] if surviving else None
        tasks.append((blk_id, source, priority, deficit, surviving))

    # Sort: CRITICAL first, then HIGH
    tasks.sort(key=lambda t: 0 if t[2] == "CRITICAL" else 1)
    return tasks, live_nodes

failed = {"dn01", "dn07"}
tasks, live = simulate_failure(failed, block_map, HEALTHY_NODES)

print(f"Failed nodes: {failed}")
print(f"Surviving nodes ({len(live)}): {sorted(live)}\n")
print(f"{'Block':<10} {'Source':<8} {'Priority':<10} "
      f"{'Deficit':>8} {'Surviving replicas'}")
print("-" * 60)
for blk, src, pri, deficit, surviving in tasks:
    print(f"{blk:<10} {str(src):<8} {pri:<10} "
          f"{deficit:>8}  {surviving}")
print(f"\n{len(tasks)} blocks require re-replication.")
\end{lstlisting}
